\chapter{Memory（记忆）}

\section{概述}

OpenClaw 记忆是\textbf{智能体工作空间中的纯 Markdown 文件}。模型只"记住"写入磁盘的内容——记忆文件是唯一的事实来源。

配置路径：\texttt{memorySearch.*}

\section{记忆文件}

默认工作空间布局使用两个记忆层：

\begin{center}
\begin{tabular}{lll}
\toprule
\textbf{文件} & \textbf{用途} & \textbf{加载时机} \\
\midrule
\texttt{memory/YYYY-MM-DD.md} & 每日日志（仅追加） & 会话开始时读取今天和昨天 \\
\texttt{MEMORY.md} & 精心整理的长期记忆 & 仅在主要私人会话中加载 \\
\bottomrule
\end{tabular}
\end{center}

位于工作空间下（默认 \texttt{\textasciitilde{}/.openclaw/workspace}）。

\section{何时写入记忆}

\begin{itemize}
  \item 决策、偏好和持久性事实 → \texttt{MEMORY.md}
  \item 日常笔记和运行上下文 → \texttt{memory/YYYY-MM-DD.md}
  \item 用户说"记住这个"时 → 立即写入
\end{itemize}

\textbf{提示}：如果你想让某些内容持久保存，请\textbf{明确要求}智能体将其写入记忆。

\section{自动记忆刷新}

当会话接近自动压缩时，OpenClaw 触发一个\textbf{静默智能体回合}，提醒模型在上下文压缩前写入持久记忆。

\begin{lstlisting}[language=json]
{
  agents: {
    defaults: {
      compaction: {
        memoryFlush: {
          enabled: true,
          softThresholdTokens: 4000
        }
      }
    }
  }
}
\end{lstlisting}

\section{向量记忆搜索}

OpenClaw 在记忆文件上构建小型向量索引，支持语义查询（即使措辞不同也能找到相关笔记）。

\subsection{默认行为}

\begin{itemize}
  \item 默认启用，监视记忆文件变更
  \item 自动选择嵌入提供商：本地 → OpenAI → Gemini
  \item 使用 sqlite-vec 加速向量搜索
\end{itemize}

\subsection{配置提供商}

\begin{lstlisting}[language=json]
{
  agents: {
    defaults: {
      memorySearch: {
        provider: "openai",
        remote: { apiKey: "YOUR_OPENAI_API_KEY" }
      }
    }
  }
}
\end{lstlisting}

支持的提供商：\texttt{openai}、\texttt{gemini}、\texttt{local}（node-llama-cpp）。

\subsection{额外记忆路径}

索引默认工作空间之外的 Markdown 文件：

\begin{lstlisting}[language=json]
{
  agents: {
    defaults: {
      memorySearch: {
        extraPaths: ["../team-docs", "/srv/shared-notes/overview.md"]
      }
    }
  }
}
\end{lstlisting}

\section{全文搜索}

除向量搜索外，OpenClaw 还支持基于 SQLite FTS5 的全文搜索，可与向量搜索结合使用以提高召回率。

\section{记忆搜索工具}

记忆搜索由活动的记忆插件提供（默认：\texttt{memory-core}）。禁用方式：

\begin{lstlisting}[language=json]
{ plugins: { slots: { memory: "none" } } }
\end{lstlisting}
